\documentclass[
		paper=a4,
	  twoside=false,
	 fontsize=10pt,
			  titlepage,
			  plainheadsepline,
			  plainfootsepline,
			  headsepline,
			  footsepline,
		index=totoc,	% fügt Inhaltsverzeichnis in Inhaltsangabe ein
	   listof=totoc,	% fügt das Sourcenverzeichnis in Inhaltsangabe ein
 bibliography=totoc,	% fügt Literaturverzeichnis in Inhaltsangabe ein
	 abstract=false,		%needed?
	titlepage=false
]{scrreprt}

\raggedright

%\parskip2ex plus1ex minus0.5ex
\setlength{\parindent}{0.05\textwidth}

\usepackage[
		 left=3cm,
		right=3cm,
		  top=2.5cm,
	   bottom=2.5cm,
		includeheadfoot
]{geometry}

\usepackage[
		pdftex,
		colorlinks=true,
		 linkcolor=black,
		 citecolor=black,
		 filecolor=black,
		  urlcolor=black
]{hyperref}
\usepackage[all]{hypcap}

\usepackage[english, ngerman]{babel}\selectlanguage{ngerman}
\usepackage[ngerman]{translator}
\usepackage[acronym, toc, numberedsection=false]{glossaries}

\usepackage[utf8]{inputenc} %Kodierung, statt latin1 auch utf8 möglich
%\usepackage{setspace} \onehalfspacing

\usepackage{lmodern}
%\usepackage{ae,aeguill}
\usepackage{ae}

\usepackage{amsfonts}
\usepackage{amsmath}
\usepackage[automark]{scrpage2}
\pagestyle{scrheadings}
%\setlength{\headheight}{1.5\baselineskip}
\clearscrheadfoot

\usepackage[round,numbers,square]{natbib}
\bibliographystyle{alphadin}

\usepackage{graphicx}
\usepackage{epsfig}
\setlength{\unitlength}{10mm} %sets unitlength for picture-environment

\usepackage{rotating}
\usepackage{longtable}
\usepackage{multirow}
\usepackage{afterpage}

% \usepackage{eurosym} % euro sign

\usepackage{color}
\definecolor{blue}{rgb}{0.1647,0,1}
\definecolor{darkred}{rgb}{0.498,0,0.33333}
\definecolor{green}{rgb}{0.247,0.498,0.3725}
\definecolor{lightgray}{rgb}{0.95,0.95,0.95}

\usepackage{listings}
\lstset{
		language=Java,									% the language for the code
	escapeinside={(*@}{@*)},							% the comments in code that are not printed
	  basicstyle=\scriptsize\ttfamily\color{black},		% the style of the code font
	keywordstyle=\bfseries\color{darkred},				% the style for keywords
	 stringstyle=\color{blue},							% the style for strings
	commentstyle=\color{green},							% the style for comments
 identifierstyle=\color{black},							% the style for identifiers
		 numbers=left,									% where to put the line numbers
	 numberstyle=\tiny,									% the style of line numbers
showstringspaces=false,									% underline spaces in strings
	  stepnumber=1,										% the step between two numbered lines
	   numbersep=10pt,									% space between numbers and code
	  captionpos=b,										% where the caption should be printed (bottom)
		   frame=single,  								% if a frame should be printed around the code
		 tabsize=4,										% how many spaces one tab contains
       aboveskip=1cm,									% the space above the displayed listing
 backgroundcolor=\color{lightgray},							% background color to fill the listings with
			breaklines
}
% \renewcommand{\lstlistlistingname}{Quellcodeverzeichnis}
%\renewcommand{\baselinestretch}{1.5}\normalsize

% Schusterjungen/Hurenkinder ausschließen
\clubpenalty = 10000
\widowpenalty = 10000
\displaywidowpenalty = 10000

\newcommand{\figref}[1]{Abbildung~\ref{fig:#1}}
\newcommand{\secref}[1]{Abschnitt~\ref{sec:#1}}
\newcommand{\tabref}[1]{Tabelle~\ref{tbl:#1}}
\newcommand{\lstref}[1]{Listing~\ref{lst:#1}}
\newcommand{\srcref}[1]{Zeile~\ref{src:#1}}
\newcommand{\chpref}[1]{Kapitel~\ref{chp:#1}}
\newcommand{\icon}{icon~Systemhaus~GmbH\ }

% \usepackage{ifthen}
% \newcommand{\vgl}[2][]{% #1 = optional, #2 = notwendig
%   \ifthenelse{\equal{#1}{}}{
%        % Autor hat kein oder leeres optionales Argument angegeben
% 		(Vgl. \cite{#2})
%   }{%
%        % Autor hat optionales Argument angegeben
% 		(Vgl. \cite[#1]{#2})
%   }
% }
% \newcommand{\vglausf}[2][]{% #1 = optional, #2 = notwendig
%   \ifthenelse{\equal{#1}{}}{
%        % Autor hat kein oder leeres optionales Argument angegeben
% 		(Vgl. dazu ausführlich \cite{#2})
%   }{%
%        % Autor hat optionales Argument angegeben
% 		(Vgl. dazu ausführlich \cite[#1]{#2})
%   }
% }

\newcommand{\code}[1]{\texttt{#1}}
\newcommand{\notiz}[1]{\marginpar{\tiny#1}}
\newcommand{\glentry}[3]{\newglossaryentry{#1}{name={#2},description={{#3}}}}

\usepackage[texcoord]{eso-pic}
%\newcommand{\startprintinglogo}{
%	\AddToShipoutPicture{
%		\parbox[t]{\paperwidth}{
%			\vspace{0.5cm}
%			\hspace{16.2cm}
%			\includegraphics[width=5.36cm,height=2.0cm]{includes/dhbw}
%		}
%	}
%}

%\renewcommand{\notiz}[1]{} %uncomment for _NOT_ showing sidenotes!
\newcommand\thema{Design und Implementierung eines Debuggers für Mikro-Assembler-Programme}

\hypersetup{
	pdftitle={Studienarbeit}, pdfauthor={Christian Rösch, icon Systemhaus GmbH},
	pdfsubject={\thema}, pdfkeywords={}
}

%\(i|c|o)head[plainpage]{otherpage} % i..links, o..rechts, c..zentriert
\ihead[Studienarbeit]{Studienarbeit}
\ohead[Christian Rösch]{Christian Rösch}
\ifoot[\rightmark]{\rightmark} 
\ofoot[\pagemark]{\pagemark} 

\loadglsentries{glossar}
\loadglsentries[\acronymtype]{abkuerzungen}
\makeglossaries

\begin{document}

	%\nocite{*} % Alle Einträge der .bib-Datei werden auch wenn unzitiert ins Literaturverzeichnis übernommen
	
	% don't write the text of these acronyms
	% \glsunset{asf}
	% \glsunset{html}

	\thispagestyle{empty}

\AddToShipoutPicture*{
	\parbox[t]{\paperwidth}{
		\vspace{2.0cm}
		\hspace{12.0cm}
		\includegraphics[height=2.5cm]{includes/dhbw.jpg}
	}
}
\parbox[t]{\paperwidth}{}
\begin{center}
	\vspace{15mm}	{\LARGE\bf\thema}\\
	\vspace{20mm}	{\Large\bf STUDIENARBEIT}\\
	\vspace{15mm}	des Studiengangs Angewandte Informatik\\der Dualen Hochschule Baden-Württemberg Stuttgart\\
	\vspace{7,5mm}	von\\
	\vspace{7,5mm}	{\bf CHRISTIAN RÖSCH}\\
	\vspace{10mm}	Oktober 2011 - Mai 2012
	\vspace{30mm}
\end{center}
\begin{longtable}{p{.51\textwidth}p{.43\textwidth}}
	\textbf{Bearbeitungszeitraum}				&  8 Monate\\
	\textbf{Matrikelnummer}						&  0487930\\
	\textbf{Kurs}								&  TIT09AIC\\
	\textbf{Ausbildungsfirma}					&  \icon - Stuttgart\\
	% \textbf{Gutachter der Ausbildungsfirma}		&  Dipl.-Inf. Steffen Huber\\
	\textbf{Gutachter der Studienakademie}	& Prof. Dr. Karl Stroetmann  \\
\end{longtable}\clearpage

	\pagenumbering{roman}
% 	\include{Erklaerung}\clearpage

%	\startprintinglogo
 	\chapter{Einleitung}
Die vorliegende Arbeit stellt einen Debugger für einen CISC-Prozessor vor. Hierbei handelt es sich um den in \cite{Stroetmann2007} vorgestellten Prozessor \emph{Mic-1}:
\begin{quote}
Die Abkürzung CISC steht hier für \emph{complex instruction set computer} und drückt aus, dass die Maschinensprache wesentlich umfangreicher ist als bei einem RISC-Prozessor.
\end{quote}

Im Gegensatz zu einem RISC-Prozessor sind die Maschinenbefehle in einem CISC-Prozessor nicht in Hardware implementiert, sondern werden durch einen Mikroprogramm-Speicher definiert. Dieser wird in Mikro-Assembler programmiert und ermöglicht es, die Maschinenbefehle flexibel zu steuern.

Der Debugger soll dazu die Entwicklung von Mikro-Assembler- und Assembler-Programmen für den CISC-Prozessor erleichtern.\clearpage
	\tableofcontents\clearpage
	\pagenumbering{arabic}

	\part{Bedienung}
\chapter{Allgemein}
\section{Parameter}
\section{Logs}
\chapter{Interaktion per Konsole}
Zunächst soll der Debugger für eine Bedienung per Kommandozeile entwickelt werden. Der Benutzer soll den Debugger starten können und anschließend über die Eingabe von verschiedenen Kommandos steuern können.

Um entsprechenden (Mikro-)Assembler-Bytecode debuggen zu können, ist es nötig den CISC-Prozessor simulieren zu können. Im Jahr 2004 wurde von Thomas Kutzer eine Studienarbeit zum Thema ``Mic-1 Simulator'' verfasst. Der Debugger soll den darin entwickelten Simulator ersetzen und um einige Funktionen ergänzen. Der Debugger soll sich daher per Kommandozeile bedienen lassen, wie der Simulator --- das heißt, er muss folgende Befehle unterstützen:
\begin{description}
\item[help] \hspace*{\fill}\\

Zeigt einen Hilfe-Text an, der die möglichen Befehle und jeweils eine kurze Beschreibung anzeigt.

\item[run] \hspace*{\fill}\\

Der Debugger simuliert die Abarbeitung des Assembler- und Mikro-Assembler-Programms bis zu einem Haltepunkt oder bis zum Programmende.

\item[step] \emph{[n]}

Der Debugger simuliert die Ausführung von \emph{n} Instruktionen des Assembler-Programms. Hierbei ist der Parameter \emph{n} eine natürliche Zahl und kann weggelassen werden; in diesem Fall wird eine Instruktion ausgeführt.

\item[reset] \hspace*{\fill}\\

Der Debugger wird in den Anfangszustand zurückgesetzt: Der Debugger verhält sich, als wäre er gerade gestartet worden.

\item[set] \emph{Register Wert}

Es wird das angegebene \emph{Register} auf den übergebenen \emph{Wert} gesetzt; dieser muss als Dezimalzahl eingegeben werden.

\item[setmemory] \emph{Adresse Wert}

Es wird das \emph{Wort} an der angegebenen \emph{Adresse} auf den übergebenen \emph{Wert} gesetzt; \emph{Adresse} und \emph{Wert} müssen als Dezimalzahl eingegeben werden. Der Speicher wird byteweise addressiert.

\item[break] \emph{Register=Wert}

Es wird ein bedingter Haltepunkt gesetzt: Sobald das angegebene \emph{Register} den Wert \emph{Wert} erhält, hält der Debugger.

\item[ls-break] \hspace*{\fill}\\

Es werden alle Haltepunkte angezeigt --- mit der jeweiligen Bedingung.

\emph{Beachte:} Dieser Befehl entspricht in etwa dem Befehl \texttt{rmbreak} aus dem Simulator von Thomas Kutzer. Da sich die Funktionalität allerdings verändert hat, wird alternativ der Befehl \texttt{ls-break} benutzt.

\item[ls-reg] \emph{[Register]}

Es wird der Inhalt von \emph{Register} angezeigt. Wird der optionale Parameter weggelassen, werden die Inhalte aller Register angezeigt.

\emph{Hinweis:} Dieser Befehl entspricht dem Befehl \texttt{show} aus dem Simulator von Thomas Kutzer.

\item[ls-mem] \emph{Start Stop}

Es wird der Inhalt des Hauptspeichers zwischen den Adressen \emph{Start} und \emph{Stop} angezeigt. Die Adressen müssen als Dezimalzahl angegeben werden und adressieren den Speicher byteweise.

\emph{Hinweis:} Dieser Befehl entspricht dem Befehl \texttt{dump} aus dem Simulator von Thomas Kutzer.

\item[trace-reg] \emph{[Register]}

Es wird das \emph{Register} zum \emph{Beobachten} registriert, dadurch wird dessen Inhalt nach jedem Simulationsschritt des Debuggers angezeigt. Wird der optionale Parameter weggelassen, werden alle Register zum \emph{Beobachten} registriert.

\emph{Hinweis:} Dieser Befehl entspricht dem Befehl \texttt{trace} aus dem Simulator von Thomas Kutzer.

\item[untrace-reg] \emph{[Register]}

Es wird das \emph{Register} vom \emph{Beobachten} deregistriert, dadurch wird dessen Inhalt nach einem Simulationsschritt des Debuggers nicht mehr angezeigt. Wird der optionale Parameter weggelassen, werden alle Register deregistriert.

\emph{Hinweis:} Dieser Befehl entspricht dem Befehl \texttt{untrace} aus dem Simulator von Thomas Kutzer.

% TODO überprüfen, ob das wirklich umsetzbar ist!
\item[trace-var] \emph{[Variable]}

Es wird die \emph{Variable} zum \emph{Beobachten} registriert, dadurch wird deren Inhalt nach jedem Simulationsschritt des Debuggers angezeigt. Wird der optionale Parameter weggelassen, werden alle Variablen zum \emph{Beobachten} registriert.

\item[untrace-reg] \emph{[Register]}

Es wird die \emph{Variable} vom \emph{Beobachten} deregistriert, dadurch wird deren Inhalt nach einem Simulationsschritt des Debuggers nicht mehr angezeigt. Wird der optionale Parameter weggelassen, werden alle Variablen deregistriert.

\item[trace-mic] \hspace*{\fill}\\

Es wird der \emph{Mikro-Assembler-Code} zum \emph{Beobachten} registriert, dadurch wird jede ausgeführte Mikro-Instruktion des Debuggers angezeigt.

\item[untrace-mic] \hspace*{\fill}\\

Es wird der \emph{Mikro-Assembler-Code} zum \emph{Beobachten} deregistriert, dadurch wird ausgeführte Mikro-Instruktionen des Debuggers nicht mehr angezeigt.

\item[trace-mac] \hspace*{\fill}\\

Es wird der \emph{Assembler-Code} zum \emph{Beobachten} registriert, dadurch wird jede ausgeführte Assembler-Instruktion des Debuggers angezeigt.

\item[untrace-mac] \hspace*{\fill}\\

Es wird der \emph{Assembler-Code} zum \emph{Beobachten} deregistriert, dadurch wird ausgeführte Assembler-Instruktionen des Debuggers nicht mehr angezeigt.

\item[exit] \hspace*{\fill}\\

Der Debugger wird beendet.
\end{description}

Zusätzlich zu den im Simulator vorhandenen Befehlen soll der Debugger folgende Befehle unterstützen:
\begin{description}
\item[micro-step] \emph{[n]}

Der Debugger simuliert die Ausführung von \emph{n} Instruktionen des Mikro-Assembler-Programms. Hierbei ist der Parameter \emph{n} eine natürliche Zahl und kann weggelassen werden; in diesem Fall wird eine Instruktion ausgeführt.

\item[rm-break] \emph{Nummer}

Es wird der Haltepunkt \emph{Nummer} entfernt. Der Parameter \emph{Nummer} ist eine natürliche Zahl und bezieht sich auf die Nummerierung, die durch \texttt{ls-break} erhalten wird. Die Haltepunkte werden anschließend neu nummeriert, so dass ein erneuter Aufruf von \texttt{ls-break} empfohlen wird.

\item[micro-break] \emph{Adresse}

Es wird ein Haltepunkt hinzugefügt, der die Simulation des Debuggers anhält, sobald der Bytecode des Mikro-Assemblers an der \emph{Adresse} ausgeführt werden soll. \emph{Adresse} wird hierbei als Dezimalzahl übergeben.

\item[macro-break] \emph{Adresse}

Es wird ein Haltepunkt hinzugefügt, der die Simulation des Debuggers anhält, sobald der Bytecode des Assemblers an der \emph{Adresse} ausgeführt werden soll. \emph{Adresse} wird hierbei als Dezimalzahl übergeben.

\item[ls-stack] \hspace*{\fill}\\

Es wird der aktuelle Inhalt des Stacks angezeigt.

\item[ls-micro-code] \emph{[Nummer1 [Nummer2]]}

Es wird der disassemblierte Programmcode des Mikro-Assemblers angezeigt. Dabei gibt es drei Varianten des Befehls:
\begin{itemize}
\item Ohne Parameter wird der komplette Mikro-Assembler-Code angezeigt.
\item Mit einem Parameter \emph{Nummer1} wird die angegebene Anzahl an Zeilen vor und nach der nächsten auszuführenden Codezeile angezeigt.
\item Mit zwei Parametern, die beide Zeilennummern darstellen, wird der Code von Zeile \emph{Nummer1} bis zur Zeile \emph{Nummer2} angezeigt. 
\end{itemize}

\item[ls-macro-code] \emph{[Nummer1 [Nummer2]]}

Dieser Befehl entspricht dem Befehl \texttt{ls-micro-code}, mit dem Unterschied, dass \texttt{ls-macro-code} sich auf den disassemblierten Assembler-Code bezieht.

\end{description}

% TODO überlegen
Unabhängig von den Befehlen soll es möglich sein gewisse Zahlen sowohl als Dezimal- als auch Hexadezimalzahl anzugeben. Möglicherweise sind auch Binär- und Oktalzahlen nötig. Es ist möglich dies über ein Parameter beim Programmstart global zu setzen, oder für jeden Befehl einzeln zu bestimmen.
\chapter{Interaktion per GUI}
\clearpage
        \part{Implementierung}
\chapter{Werkzeuge}
Mit diesem Teil der Arbeit möchte ich besonders eine Frage beantworten: Wie sind die bisher vorgestellten Funktionen des Debuggers implementiert? Mein Ziel ist, die Implementierung und Struktur des Debuggers so nachvollziehbar zu beschreiben, dass Sie in der Lage sind den Debugger weiterzuentwickeln.

Der Debugger soll auch nach Ende dieser Arbeit weiterentwickelt werden können. Daher stelle ich in diesem Kapitel vor, welche Werkzeuge ich zur Entwicklung des Debuggers nutze und wie Sie diese zur Mitarbeit nutzen können. Dies ist nicht als Bedienungsanleitung zu verstehen. Ich bin der Ansicht, dass es das Verständnis der Implementierung und meiner Arbeit steigert, wenn geklärt ist, welche Werkzeuge ich wie nutze.

\section{Versionskontrollsystem}
Das Projekt ist bei github % TODO Referenz
gehostet und ist somit mit \gls{git} versioniert.

\gls{git} % Quelle: http://git-scm.com/about 04.12.2011 14:00
ist:
\begin{description}
\item[effizient] Auch bei großen Projekten zeigen Vergleiche, %TODO Beleg
dass \gls{git} insgesamt schneller ist, als andere Versionskontrollsysteme.
\item[verteilt] Im Gegensatz zu beispielsweise CVS % TODO CVS
ist \gls{git} ein verteiltes Versionskontrollsystem: Jeder Entwickler erhält ein lokales Repository %TODO repository
und kann damit arbeiten. Es wird kein zentraler Server benötigt, alle Funktionen des Systems und Commit-Informationen stehen lokal zur Verfügung.
\item[sicher] %TODO nochmal genau nachschauen
Bei \gls{git} werden viele Hash-Summen gebildet und sobald ein Commit veröffentlich wurde, ist es nicht möglich diesen zu manipulieren, da dies die Hash-Summe %TODo welche genau?
verändern würde.
\end{description}

Ich habe \gls{git} gewählt, da dadurch später leicht forks %TODO fork
gebildet werden können und das Projekt für andere Arbeiten weitergeführt werden kann. Das dezentrale Entwickeln mit \gls{git} ermöglicht es zusätzlich, dass später viele Entwickler das Projekt weiterentwickeln.

Ich möchte nun kurz erläutern, wie man mit \gls{git} den Quelltext der Arbeit auscheckt. Zusätzlich möchte ich erläutern, wie Sie \gls{git} unterstützt, wenn Sie das Projekt weiterentwickeln möchten. Dadurch möchte ich keine umfassende Bedienungsanleitung für \gls{git} geben, sondern Ihnen die alltägliche Arbeit mit dem Werkzeug erklären.

\subsection{Code auschecken}
\label{sec:git-checkout}
Die Adresse des öffentlichen Repositories %TODO genitiv repository
lautet:
\code{https://github.com/croesch/micro-debug.git}

Das Auschecken des Repositories %TODO hier auch
erzeugt ein lokales Repository im Verzeichnis \emph{micro-debug}; mit folgenden Befehlen können Sie das Repository auschecken und anschließend in das erzeugte Verzeichnis wechseln:
\begin{verbatim}
$ git clone git://github.com/croesch/micro-debug.git
$ cd micro-debug
\end{verbatim}

Im Gegensatz zu SVN %TODO svn
erstellt \gls{git} nur im Wurzelverzeichnis eines Projektes einen Ordner \emph{.git}, dieser enthält alle Informationen über das Repository. Direkt nach dem Klonen eines anderen Repositories enthält es beispielsweise einen Verweis (\emph{origin}) auf dieses Repository, um Aktualisierungen zu erhalten.

Im Laufe der Zeit wird das öffentliche Repository aktuellere Commits enthalten, als Ihr lokales Repository. Aktualisiert man mit \gls{git} das Repository,
\begin{verbatim}
$ git pull origin
\end{verbatim}
so werden die neuen Commits heruntergeladen und direkt in den aktuell ausgecheckten Branch gemergt. %TODO denglisch

Dieses Vorgehen kann manchmal unerwünscht sein, so dass man zunächst das Repository aktualisieren möchte und später selbst mergen möchte:
\begin{verbatim}
$ git fetch origin
$ git merge origin/master
\end{verbatim}
Mit \emph{fetch} werden die Aktualisierungen zunächst nur im \emph{.git}-Verzeichnis gespeichert. Erst \emph{merge} verändert die lokalen Dateien im aktuellen Branch.

\subsection{Code beitragen}
Hat man ein lokales Repository ausgecheckt und möchte Codeänderungen dem Projekt beitragen, gibt es zwei Möglichkeiten: Man benutzt den Standard-Branch oder so genannte Feature-Branches.

Den Standard-Branch zu nutzen, empfehle ich nur für die Entwicklung als Einzelperson. Entwickelt man alleine an einem Projekt, kann man veränderte Dateien wie folgt committen:
\begin{verbatim}
$ git add relativer-pfad-zur-datei
$ git commit -m "Commit-Nachricht"
\end{verbatim}

Für mehrköpfige Entwicklungsteams empfehle ich Feature-Branches zu nutzen, da \gls{git} gute branch-/merge-Qualitäten aufweist. Beispielsweise können Sie so selbst Code zu dem Projekt beitragen, indem Sie folgende Schritte durchführen:
\begin{enumerate}
\item Erstellen Sie unter \code{https://github.com/croesch/micro-debug/issues} ein Thema, das Ihre Entwicklung beschreibt.

\item Erstellen Sie einen lokalen Branch zu dem Thema, wobei die Nummer die Nummer des Themas ist:
\begin{verbatim}
$ git checkout -b 100-themen-titel
\end{verbatim}

\item Entwickeln Sie den Code weiter.

\item Committen Sie die Änderungen:
\begin{verbatim}
$ git add .
$ git commit -m 'Commit-Nachricht'
\end{verbatim}

\item Laden Sie den Stand des öffentlichen Repositories:
\begin{verbatim}
$ git fetch origin
\end{verbatim}

\item Aktualisieren Sie den lokalen master-Branch:
\begin{verbatim}
$ git checkout master
$ git pull origin master
\end{verbatim}

\item Wiederholen Sie die Schritte 2-6, bis Ihre Entwicklung abgeschlossen ist

\item Aktualisieren Sie den Punkt von dem Ihr Branch abgeht:
\begin{verbatim}
$ git checkout 100-themen-titel
$ git rebase master
\end{verbatim}

\item Stellen Sie Ihren Branch öffentlich zur Verfügung.
\begin{verbatim}
$ git push origin 100-themen-titel
\end{verbatim}
\end{enumerate}

Da Sie für dieses Projekt auf github keine Schreibberechtigung haben, ist der letzte Punkt nicht so einfach möglich. Daher wählt man besonders auf gitHub einen leicht anderen Weg, den ich nun kurz vorstellen möchte.

\subsection{Unterstützung durch gitHub}
Die Grundidee ist, ein beliebiges Repository auf github zu klonen, so dass dieser Klon wiederum öffentlich zugänglich ist. Sie klonen dann den Klon von github auf Ihren lokalen Rechner, um daran zu entwickeln. Anschließend können Sie Ihr Ergebnis durch Ihren github-Klon veröffentlichen, so dass diese Änderung wieder in das Originalprojekt einfließen kann. Diesen Prozess möchte ich nun erläutern.

\begin{enumerate}
\item Registrieren Sie sich auf github, um dort Projekte veröffentlichen zu können. Legen Sie ssh-Schlüsselpaare an und laden Sie Ihren öffentlichen Schlüssel auf Ihr Benutzerkonto, dass Sie sich später bei Schreiboperationen authentifizieren können.

\item Klicken Sie auf der Projektseite \code{https://github.com/croesch/micro-debug} den \emph{fork}-Button. Dadurch klont github das Projekt auf Ihr Benutzerkonto, wo es nun unter \code{https://github.com/ihr-name/micro-debug} verfügbar ist.

\item Klonen Sie Ihre Repository nun lokal, wie in \secref{git-checkout} beschrieben:
\begin{verbatim}
$ git clone git://github.com/ihr-name/micro-debug.git
$ cd micro-debug
\end{verbatim}

\item Da Ihr Repository unter \emph{origin} nun Ihren github-Klon referenziert, setzen Sie eine Referenz auf das Hauptprojekt:
\begin{verbatim}
$ git remote add original-projekt git://github.com/croesch/micro-debug.git
\end{verbatim}

\item Erstellen Sie unter \code{https://github.com/croesch/micro-debug/issues} ein Thema, das Ihre Entwicklung beschreibt.

\item Erstellen Sie einen lokalen Branch zu dem Thema, wobei die Nummer die Nummer des Themas ist:
\begin{verbatim}
$ git checkout -b 100-themen-titel
\end{verbatim}

\item Entwickeln Sie den Code weiter.

\item Committen Sie die Änderungen:
\begin{verbatim}
$ git add .
$ git commit -m 'Commit-Nachricht'
\end{verbatim}

\item Laden Sie den Stand des öffentlichen Repositories:
\begin{verbatim}
$ git fetch original-projekt
\end{verbatim}

\item Aktualisieren Sie den lokalen master-Branch:
\begin{verbatim}
$ git checkout master
$ git pull original-projekt master
\end{verbatim}

\item Wiederholen Sie die Schritte 2-6, bis Ihre Entwicklung abgeschlossen ist

\item Aktualisieren Sie den Punkt von dem Ihr Branch abgeht:
\begin{verbatim}
$ git checkout 100-themen-titel
$ git rebase master
\end{verbatim}

\item Veröffentlichen Sie Ihren fertigen Branch über Ihren github-Klon:
\begin{verbatim}
$ git push origin 100-themen-titel
\end{verbatim}

\item Klicken Sie auf Ihrer Projekt-Seite auf \emph{Pull-Request}. Anschließend kann Ihr neuer Branch beispielsweise von mir in das Originalprojekt eingepflegt werden. Dann können Sie und alle anderen, die das Projekt geklont haben, die Änderungen über \emph{fetch} in die geklonten Repositories übernehmen.
\end{enumerate}

Sie haben nun die Möglichkeit den aktuellen Code lokal zu bearbeiten und Änderungen zum Projekt beizutragen. Im Folgenden möchte ich das Build-Management-Werkzeug erläutern, dass ich zum Entwickeln nutze: \gls{mvn}.

\section{Build-Management}
\gls{mvn} ist ein Werkzeug zum Bauen von Java-Anwendungen. Meines Erachtens gibt es keine signifikanten Gründe für oder gegen die Nutzung von \gls{mvn}; stattdessen hätte ich auch \gls{ant} oder \gls{make} nutzen können. In diesem Abschnitt möchte ich die grundlegenden Befehle für die Arbeit mit \gls{mvn} erläutern.

Ich verwende \gls{mvn}, da hier viel implizites Wissen eingesetzt wird; dadurch wird der Konfigurationsaufwand geringer -- solange man den Konventionen der \gls{mvn}-Entwickler folgt. Zu dem impliziten Wissen gehört unter anderem:
\begin{itemize}
\item Java-Klassen befinden sich in \code{src/main/java/}
\item Java Test-Klassen befinden sich in \code{src/test/java/}
\item Kompilate und von \gls{mvn} generierte Klassen befinden sich in \code{target/}
\item Die explizite Konfiguration steht in der Datei \code{pom.xml}
\end{itemize}

\subsection{Code kompilieren}
Mit dem folgenden Befehl können Sie den Programmcode kompilieren; Testklassen werden dadurch weder kompiliert noch ausgeführt:
\begin{verbatim}
$ mvn compile
\end{verbatim}

\gls{mvn} erkennt verschiedene Abhängigkeiten. Das führt dazu, dass beim ersten Ausführen zunächst erstmal alle benötigten Bibliotheken heruntergeladen werden. Zusätzlich erkennt \gls{mvn} auch, ob für ein bestimmtes Ziel die Klassen zunächst kompiliert werden müssen.

Beispielsweise werden zum Packetieren einer Anwendung zunächst alle Klassen kompiliert, anschließend die Tests ausgeführt, bevor die Anwendung packetiert werden kann; dies geschieht mit:
\begin{verbatim}
$ mvn package
\end{verbatim}

\subsection{Tests ausführen}
Wenn Sie alle Tests ausführen möchten, können Sie dies mit folgendem Befehl erreichen:
\begin{verbatim}
$ mvn test
\end{verbatim}

Je nachdem, welche Metriken Sie per Plugins noch mit \gls{mvn} ausführen, kann es nötig sein die von \gls{mvn} generierten Dateien zu löschen. Beispielsweise bei Code-Coverage-Werkzeugen ist dies nötig; alle von \gls{mvn} generierten Dateien -- das Verzeichnis \code{target} -- entfernen Sie mit:
\begin{verbatim}
$ mvn clean
\end{verbatim}

\subsection{Abhängigkeiten ändern}
Die Date \code{pom.xml} enthält die explizite Konfiguration eines Projektes. Darin enthalten sind unter anderem der Name und die Version des Projekts, aber auch abhängige Projekte.

Beispielsweise nutzt diese Arbeit \gls{fest} als Test-Werkzeug, in der \code{pom.xml} steht dazu Folgendes:
\begin{verbatim}
<dependency>
    <groupId>org.easytesting</groupId>
    <artifactId>fest-swing</artifactId>
    <version>1.2.1</version>
    <scope>test</scope>
</dependency>
\end{verbatim}

Dort können Sie nun die Version des Werkzeugs ändern, die genutzt werden soll; oder den Bereich von \emph{test} auf \emph{runtime} erweitern. Anschließend empfehle ich die Dateien zu entfernen, die mit \gls{mvn} zuvor generiert wurden und \gls{mvn} auszuführen.

\section{Entwicklungsumgebung}
\subsection{Integration von git und maven}
\subsection{Navigation im Code}
\subsection{Finden von Referenzen}
\chapter{Prozessor Mic1}
\section{ALU}
\section{MPC-Berechnung}
\section{Mic1-Instruktionen}
\section{Register}
\section{Hauptspeicher}
\section{Zusammensetzung der Komponenten}
\chapter{Interaktion per Konsole}
\chapter{Interaktion per GUI}\clearpage

	\appendix
	\pagenumbering{Roman}
	\printglossary\clearpage
	\printglossary[type=\acronymtype, title={Abkürzungsverzeichnis}, toctitle={Abkürzungsverzeichnis}]\clearpage
	\listoffigures\clearpage
	%\listoftables\clearpage
	\lstlistoflistings\clearpage
	\bibliography{includes/literatur}\clearpage
	\chapter*{Erklärung}\addcontentsline{toc}{chapter}{Erklärung}
\markright{Erklärung}
\begin{flushleft}
gemäß \S 5 (2) der "`Studien- und Prüfungsordnung DHBW Technik"' vom 18. Mai 2009.\\

Ich versichere hiermit, dass ich die vorliegende Arbeit mit dem Thema\\
\vspace{10mm}\thema\\\vspace{10mm}
selbstständig verfasst und keine anderen als die angegebenen Quellen und Hilfsmittel verwendet habe.\\
\vspace{30mm}
\begin{tabular}{lp{15mm}p{70mm}}
	Stuttgart, \today && \\\cline{1-1}\cline{3-3}
	{\footnotesize Ort, Datum} && {\footnotesize Christian Rösch} \\
\end{tabular}
\end{flushleft}
\clearpage

\end{document}
